\documentclass[11pt,a4paper]{article}
\usepackage[margin=2.5cm]{geometry}
\usepackage{amsmath,amssymb,graphicx,booktabs,hyperref,caption}
\usepackage{times}
\newcommand{\Hmat}{\mathbf{H}}
\begin{document}

\section*{Supplementary Note 5: Reproducibility Audit of the G-VAE Directional-Locking Positive Control}
\label{sup:opposition_audit}

The G-VAE spectral-shaping objective (a Hutchinson Jacobian-vector-product norm penalty plus an opposition gene-weight curriculum) was originally reported to produce directional locking in 21/21 seeds across three tissues, including 15/15 hematopoietic seeds with $C_{\text{head}}=77.6\%$ and all 64 vanilla directions opposed in every seed. This note documents why that positive control is withdrawn.

\subsection*{Recomputation from the surviving checkpoints}

Recomputing the opposition-budget metrics from the surviving hematopoietic checkpoints (all 15 paired vanilla/G-VAE seeds in \texttt{gvae\_hema/} and \texttt{gvae\_hema\_scan/}) with the published Ridge-readout protocol gives \textbf{0/15 seeds with $C_{\text{head}}>50\%$} ($C_{\text{head}}$ range $-67\%$ to $+2.4\%$; opposed directions 4--10 of 64, never 64). The vanilla trace matches the training-time \texttt{results.json} exactly, confirming the recomputation protocol is faithful. The training-time logs are consistent: all 15 seeds are recorded as \texttt{status=MARGINAL} with \texttt{predicted\_lock=false}.

The analysis file that reported the positive result (\texttt{results/OPPOSITION\_BUDGET\_ANALYSIS.json}; hema 15/15, $C_{\text{head}}=77.6\%$, $n_{\mathrm{opp}}=64/64$) is not derivable from any surviving checkpoint or preprocessing variant (raw, log1p, or zscore over 1{,}000--5{,}000 cells), and its generator script is absent from the repository and from git history.

\subsection*{Fresh reproduction}

A fresh three-seed reproduction (seeds 42/123/456) of the released hematopoietic recipe locks in only 1/3 seeds (s456 $C=96.1\%$; s123 $C=83.0\%$; s42 $C=46.1\%$), and only when checkpoints are selected from the first 160 epochs. Locking therefore behaves as an early-training transient preserved by checkpoint selection rather than a robust converged state. The locking is also protocol-dependent: re-auditing the same three reproduction checkpoints under the unified audit protocol (shared R14 iter-4 stagedata; StandardScaler Ridge readout, $\alpha=1.0$; G-VAE 15-seed vanilla reference $\mathrm{tr}(\Hmat_V)=4080.8$) gives $\kappa=1.02$, $0.58$, $4.78$ for seeds 42/123/456---none below the locking threshold $\kappa<0.5$ (\texttt{results/positive\_control\_repro\_audit.json}). The training-time $C$ values (46--96\%) were computed against a v13-specific vanilla baseline and checkpoint-selection protocol and do not transfer to the unified protocol; no trained model in this project, including the reproduction checkpoints, satisfies the audit's locking criterion. In contrast, a constructive sensitivity control (Supplementary Fig.~S5) shows that the audit metric does report locking when a latent is genuinely rank-collapsed (zeroing all but the top-$k$ Ridge dimensions of a real vanilla encoder gives $\kappa=0.15$--$0.44$ for $k=1$--$16$), so the 0/33 null cannot be attributed to an insensitive metric.

\subsection*{Neural-crest and lung figures}

The neural-crest (reported $C=56$--$69\%$) and lung-fibrosis (reported $C=96.7\pm0.3\%$) G-VAE figures were produced with the $\kappa$-$C$ penalty term, which is non-differentiable as implemented (Supplementary Note~4); these figures remain unverified and should be re-checked with a differentiable spectral objective before use.

\subsection*{Conclusion}

None of the G-VAE positive-control figures survives verification. Spectral locking remains architecturally plausible but is not demonstrated by the surviving evidence; the robust empirical claim of this work is the 0/33 prior-injection null, and the 21/21 contrast is withdrawn. Recomputations are archived in \texttt{results/opp\_recompute\_from\_ckpt.json} and described in \texttt{results/OPPOSITION\_BUDGET\_REPRODUCIBILITY\_AUDIT.md}.

\vspace{10pt}
\noindent\rule{\textwidth}{0.4pt}
\footnotesize\textbf{Source identity.} This note was compiled from \texttt{supplementary/supp\_note5\_opposition\_repro.tex} in the paper1 repository. SHA256 of that source file with this paragraph removed: \texttt{3b5cd3e0070d84f62936a62f8e7a2b193f9d07414dd8c95fb8d4b00df82ceef7}. Readers can recompute the digest to confirm which revision produced this PDF; the source is distributed with the code release.

\end{document}
